\section*{Step 287: abc Conjecture Classification}

Let $\operatorname{rad}(n)=\prod_{p\mid n}p$.  The abc conjecture asserts that for every
$\varepsilon>0$ there are only finitely many coprime positive integer triples
$(a,b,c)$ with $a+b=c$ and
\[
  c>\operatorname{rad}(abc)^{1+\varepsilon}.
\]
Equivalently, for every $\varepsilon>0$ there is a constant $K(\varepsilon)$ such that
\[
  c\le K(\varepsilon)\operatorname{rad}(abc)^{1+\varepsilon}.
\]

Define the residual
\[
\Xi_{\mathrm{abc}}(K,\varepsilon)=
\#\{(a,b,c): a+b=c,\ (a,b,c)=1,\ c>K\operatorname{rad}(abc)^{1+\varepsilon}\}.
\]
Closure is finiteness of the exceptional locus for every $\varepsilon>0$.

This residual is not an $L$-function zero-location, correlation, growth, moment, or
zero-density residual.  The abstract framework findings can audit proof strategies, but
they do not classify abc as a native instance.  The step therefore records a candidate
eleventh family: integer-Diophantine radical-height residuals, with abc as the primary
instance and Hall/Pillai/Erdos--Straus style residuals as adjacent analogues.

\[
\boxed{\texttt{V\_abc\_11th\_finding}}
\]
